\documentclass[11pt,a4paper]{article}
\usepackage{fontspec}
\usepackage{geometry}
\usepackage{listings}
\usepackage{xcolor}
\usepackage{hyperref}
\usepackage{titlesec}
\usepackage{enumitem}
\usepackage{booktabs}
\usepackage{parskip}

\geometry{margin=2.5cm}

\definecolor{codebg}{RGB}{245,245,245}
\definecolor{codestr}{RGB}{0,128,0}
\definecolor{codekw}{RGB}{0,0,200}
\definecolor{codecom}{RGB}{128,128,128}
\definecolor{purple}{RGB}{124,77,255}

\lstdefinelanguage{Kotlin}{
  keywords={fun, val, var, class, object, if, else, when, return, true, false,
            null, actual, expect, try, catch, override, private, internal},
  keywordstyle=\color{codekw}\bfseries,
  comment=[l]{//},
  commentstyle=\color{codecom}\itshape,
  stringstyle=\color{codestr},
  morestring=[b]",
  morestring=[b]',
  basicstyle=\ttfamily\small,
  backgroundcolor=\color{codebg},
  frame=single,
  framerule=0.4pt,
  rulecolor=\color{gray!40},
  breaklines=true,
  showstringspaces=false,
  tabsize=4
}
\lstset{language=Kotlin}

\hypersetup{
  colorlinks=true,
  linkcolor=purple,
  urlcolor=purple
}

\title{\textbf{Integracja \texttt{symja\_android\_library} z RUTEdu}\\
  \large Dokumentacja architektury i przewodnik programisty}
\author{RUTEdu Team}
\date{\today}

\begin{document}
\maketitle
\tableofcontents
\newpage

% -------------------------------------------------------------------------
\section{Wprowadzenie}

RUTEdu integruje \textbf{matheclipse-core} (jądro projektu
\href{https://github.com/axkr/symja_android_library}{symja\_android\_library}) jako silnik
algebry symbolicznej (CAS — \emph{Computer Algebra System}).

\subsection{Co umożliwia integracja}
\begin{itemize}
  \item Sprawdzanie odpowiedzi algebraicznych przez równoważność matematyczną
        (użytkownik może wpisać \texttt{2*x+6*x} zamiast \texttt{8*x} i dostać zaliczenie).
  \item Pytania o pochodne, całki, upraszczanie wyrażeń i rozwiązywanie równań.
  \item Wskazówki krok po kroku generowane na podstawie reguł rachunkowych.
\end{itemize}

\subsection{Ograniczenie platformy}
Biblioteka jest oparta na JVM — działa \textbf{tylko na Androidzie}.
Na iOS \texttt{MathEngine} to stub zwracający \texttt{null}/\texttt{false},
a \texttt{AlgebraQuestionGenerator} automatycznie serwuje pytania \texttt{SelectFromList}
z predefiniowanymi odpowiedziami.

% -------------------------------------------------------------------------
\section{Architektura: \texttt{expect/actual MathEngine}}

\subsection{Diagram}

\begin{center}
\begin{tabular}{lll}
\toprule
\textbf{Moduł} & \textbf{Plik} & \textbf{Implementacja} \\
\midrule
\texttt{commonMain} & \texttt{math/MathEngine.kt}
  & \texttt{expect class} + \texttt{expect val mathEngineAvailable} \\
\texttt{androidMain} & \texttt{math/MathEngine.android.kt}
  & \texttt{actual} wrappujący \texttt{ExprEvaluator} \\
\texttt{iosMain} & \texttt{math/MathEngine.ios.kt}
  & \texttt{actual} stub — wszystkie metody $\to$ \texttt{null} \\
\bottomrule
\end{tabular}
\end{center}

\subsection{Interfejs \texttt{commonMain}}

\begin{lstlisting}
expect val mathEngineAvailable: Boolean  // true na Android, false na iOS

expect class MathEngine() {
    val isAvailable: Boolean

    // Zwraca true jesli wyrazenia sa matematycznie rownowazne
    fun areEquivalent(userExpr: String, correctExpr: String): Boolean?

    // Uproszczona postac wyrazenia (Simplify)
    fun simplify(expr: String): String?

    // Pochodna wzgledem zmiennej (D[expr, variable])
    fun differentiate(expr: String, variable: String = "x"): String?

    // Calka nieoznaczona (Integrate[expr, variable])
    fun integrate(expr: String, variable: String = "x"): String?

    // Rozwiazanie rownania (Solve[equation, variable])
    fun solve(equation: String, variable: String = "x"): String?
}
\end{lstlisting}

\subsection{Implementacja Android}

\begin{lstlisting}
actual class MathEngine actual constructor() {
    private val evaluator: ExprEvaluator? = try {
        ExprEvaluator(false, 100)
    } catch (e: Throwable) { null }

    actual val isAvailable: Boolean get() = evaluator != null

    actual fun areEquivalent(userExpr: String, correctExpr: String): Boolean? {
        val eval = evaluator ?: return null
        return try {
            eval.eval("Simplify(($userExpr)-($correctExpr))").toString() == "0"
        } catch (e: Throwable) { null }
    }
    // ... pozostale metody analogicznie
}
\end{lstlisting}

% -------------------------------------------------------------------------
\section{Zasady bezpieczeństwa wątkowego}

\begin{itemize}[leftmargin=*]
  \item \textbf{\texttt{ExprEvaluator} NIE jest thread-safe.}
        Każde wywołanie twórz nową instancję \texttt{MathEngine()} — nigdy nie
        przechowuj jako singleton.
  \item \textbf{Inicjalizacja trwa 500--2000~ms} (pierwsze ładowanie klas JVM).
        Zawsze konstruuj na \texttt{Dispatchers.Default}, nigdy na głównym wątku.
  \item Użyj \texttt{mathEngineAvailable} do sprawdzenia platformy bez tworzenia
        instancji.
\end{itemize}

Wzorzec użycia w composable (z \texttt{ExpressionTypeAnswerContent}):

\begin{lstlisting}
val scope = rememberCoroutineScope()

fun check() {
    scope.launch(Dispatchers.Default) {    // wątek tła
        val engine = MathEngine()          // nowa instancja per wywołanie
        val ok = engine.areEquivalent(input.trim(), question.correctExpr)
        withContext(Dispatchers.Main) {
            if (ok == true) onCorrect() else onWrong()
        }
    }
}
\end{lstlisting}

% -------------------------------------------------------------------------
\section{Nowy typ pytania: \texttt{ExpressionTypeAnswer}}

\begin{lstlisting}
data class ExpressionTypeAnswer(
    override val id: Int,
    val prompt: String,
    val correctExpr: String,       // notacja Symja, np. "6*x+2"
    val displayCorrect: String = correctExpr,  // czytelna forma do feedbacku
    val inlineHint: String? = "Użyj * dla mnożenia, ^ dla potęgi",
    val hint: Hint = Hint("")
) : Question(id)
\end{lstlisting}

Composable renderujący: \texttt{ExpressionTypeAnswerContent.kt}.
Branch w \texttt{LessonGameScreen.kt}:

\begin{lstlisting}
is Question.ExpressionTypeAnswer -> ExpressionTypeAnswerContent(
    question = question,
    accentColor = accentColor,
    bottomPadding = bottomPadding,
    onCorrect = onAnsweredCorrectly,
    onWrong = onWrongAnswer
)
\end{lstlisting}

% -------------------------------------------------------------------------
\section{AlgebraQuestionGenerator}

\subsection{Schemat generatora}

\begin{lstlisting}
object AlgebraQuestionGenerator {

    fun generateFor(lessonId: String, seed: Long,
                    excludeIds: Set<Int> = emptySet()): List<Question> {
        val all = if (mathEngineAvailable) {
            when (lessonId) {
                "algebra_1_1" -> algebra_1_1_android()  // ExpressionTypeAnswer
                "algebra_1_2" -> algebra_1_2_android()
                // ...
            }
        } else {
            when (lessonId) {
                "algebra_1_1" -> algebra_1_1_ios()      // SelectFromList fallback
                "algebra_1_2" -> algebra_1_2_ios()
                // ...
            }
        }
        return all.filter { it.id !in excludeIds }
    }
}
\end{lstlisting}

\subsection{Lekcje algebry}

\begin{center}
\begin{tabular}{llll}
\toprule
\textbf{ID lekcji} & \textbf{Temat} & \textbf{Android} & \textbf{iOS fallback} \\
\midrule
\texttt{algebra\_1\_1} & Upraszczanie wyrażeń & \texttt{ExpressionTypeAnswer} & \texttt{SelectFromList} \\
\texttt{algebra\_1\_2} & Pochodne & \texttt{ExpressionTypeAnswer} & \texttt{SelectFromList} \\
\texttt{algebra\_1\_3} & Całki nieoznaczone & \texttt{ExpressionTypeAnswer} & \texttt{SelectFromList} \\
\texttt{algebra\_2\_1} & Równania liniowe & \texttt{ExpressionTypeAnswer} & \texttt{SelectFromList} \\
\texttt{algebra\_2\_2} & Równania kwadratowe & \texttt{ExpressionTypeAnswer} & \texttt{SelectFromList} \\
\bottomrule
\end{tabular}
\end{center}

\subsection{Routing w QuestionBank}

\begin{lstlisting}
fun questionsFor(lessonId: String, seed: Long = 0L,
                 excludeIds: Set<Int> = emptySet()): List<Question> =
    when {
        lessonId.startsWith("chemia_") ->
            ChemistryQuestionGenerator.generateFor(lessonId, seed, excludeIds)
        lessonId.startsWith("algebra_") ->
            AlgebraQuestionGenerator.generateFor(lessonId, seed, excludeIds)
        else ->
            banks[lessonId] ?: emptyList()
    }
\end{lstlisting}

% -------------------------------------------------------------------------
\section{Dodawanie nowej lekcji algebry}

\begin{enumerate}
  \item Dodaj \texttt{Lesson} do \texttt{SubjectRepository.kt} z ID
        zaczynającym się od \texttt{"algebra\_"}.
  \item W \texttt{AlgebraQuestionGenerator} dopisz dwie funkcje prywatne:
        \texttt{algebra\_X\_Y\_android()} i \texttt{algebra\_X\_Y\_ios()}.
  \item Zarejestruj obie w bloku \texttt{when} wewnątrz \texttt{generateFor()}.
\end{enumerate}

\textbf{Przykład — nowa lekcja \texttt{algebra\_1\_4} (Granice):}

\begin{lstlisting}
// SubjectRepository.kt
Lesson(
    id = "algebra_1_4",
    name = "Granice funkcji",
    description = "Oblicz granicę",
    progress = 0.0f, isLocked = false,
    color = Color(0xFF7C4DFF),
    icon = Icons.Default.Functions
)

// AlgebraQuestionGenerator.kt
private fun algebra_1_4_android(): List<ExpressionTypeAnswer> = listOf(
    ExpressionTypeAnswer(0,
        prompt = "Oblicz: lim(x→2) (x² - 4)/(x - 2)",
        correctExpr = "4", displayCorrect = "4",
        hint = Hint("Usuń wspólny czynnik (x-2) z licznika.",
            steps = listOf("(x²-4)/(x-2) = (x+2)(x-2)/(x-2)", "= x+2", "lim(x→2) = 4")))
)

private fun algebra_1_4_ios(): List<SelectFromList> = listOf(
    SelectFromList(0, "lim(x→2) (x²-4)/(x-2) = ?",
        listOf("4", "0", "2", "nieskończoność"), setOf(0))
)
\end{lstlisting}

% -------------------------------------------------------------------------
\section{Format wyrażeń akceptowanych przez Symja}

\begin{center}
\begin{tabular}{lll}
\toprule
\textbf{Operacja} & \textbf{Notacja Symja} & \textbf{Przykład} \\
\midrule
Mnożenie & \texttt{*} & \texttt{3*x^2} \\
Potęga & \texttt{\textasciicircum} & \texttt{x\textasciicircum 3} \\
Pierwiastek & \texttt{Sqrt[]} & \texttt{Sqrt[x]} \\
Ułamek & \texttt{/} & \texttt{x\textasciicircum 2/2} \\
Równanie & \texttt{==} & \texttt{2*x+4==0} \\
Sinus/Cosinus & \texttt{Sin[], Cos[]} & \texttt{Sin[x]} \\
\bottomrule
\end{tabular}
\end{center}

\textbf{Dla \texttt{areEquivalent}}: Symja sprawdza \texttt{Simplify(userExpr - correctExpr) == 0}.
Dzięki temu \texttt{6x+2}, \texttt{2+6*x} i \texttt{6*x+2} są traktowane jako równoważne.

% -------------------------------------------------------------------------
\section{Zależność Gradle i rozmiar APK}

\subsection{Dodawanie zależności}

W \texttt{composeApp/build.gradle.kts}:
\begin{lstlisting}
androidMain.dependencies {
    implementation(libs.matheclipse.core)
}
\end{lstlisting}

W \texttt{gradle/libs.versions.toml}:
\begin{lstlisting}
[versions]
matheclipse = "3.1.0"

[libraries]
matheclipse-core = { module = "org.matheclipse:matheclipse-core",
                     version.ref = "matheclipse" }
\end{lstlisting}

\subsection{Rozmiar APK}

Symja + zależności transytywne dodają ok. \textbf{8--15~MB} do APK.
R8/ProGuard są włączone w buildzie release (\texttt{isMinifyEnabled = true}).
Reguły keep-ów w \texttt{androidApp/proguard-rules.pro} chronią klasy
\texttt{org.matheclipse.**}, \texttt{org.antlr.**} i \texttt{org.hipparchus.**}.

\end{document}
